\chapter{Prelucrarea fișierelor audio}
\label{section:extraction}
Odată ce toate fișierele audio au fost descărcate, separate în voce și instrumental iar fișierele ce conțin vocea au fost curățate prin
procesul descris în capitolul anterior, a fost nevoie de un pipeline de extragere a datelor din fișierele audio. Pipeline-ul are scopul de a
face procesul de extragere
destul de flexibil încât să permită rularea de experimente pentru segmente de diferite dimensiuni și diferite tipuri de reprezentări ale
informațiilor audio. Pașii implementați pot fi observați în figura \ref{fig:f7}

\begin{figure}[h]
\includegraphics[width=0.80\textwidth]{feature_extraction.png}
\centering
\captionsetup{justification=centering,font=small,labelfont=small}
\caption[something]{Procesul de extragere al informațiilor din fișierele audio (Waveform preluat din \textit{Cloudinary Blog} \protect\footnotemark[1])}
\label{fig:f7}
\end{figure}
\footnotetext[1]{\url{https://cloudinary.com/blog/how_to_generate_waveform_images_from_audio_files}}
Pipeline-ul primește ca input directorul în care se află fișierele audio, directorul sau fișierul în care vor fi salvate informațiile extrase,
în funcție de forma sub care dorim să le salvăm, o listă de instanțe ale clasei \textit{Feature Function}, o instanță a clasei \textit{Segment Splitter} și,
opțional, un procent în cazul în care se dorește procesarea unui subset al setului de date, opțiune utilă în procesul de testare al
funcționalităților pe diferite mașini.

Clasa \textit{Segment Splitter} se ocupă de împărțirea semnalului audio în mai multe segmente, de dimensiuni fixe. La crearea unei instanțe pot fi
specificate numărul de secunde pe care îl vor avea segmentele și o listă de offset-uri de la care să înceapă separarea. De exemplu, după cum
am exemplificat și în figura \ref{fig:f7} dacă un offset nu este specificat, în urma separării vor rezulta doar segmentele reprezentate cu
roșu, iar dacă specificăm un offset, vor rezulta și segmentele reprezentate cu verde. Astfel, se asigură faptul că toate datele de input
pentru rețea au aceeași dimensiune prin prisma faptului că segmentele rezultate au lungimi egale. De asemenea, specificând unul sau mai multe offset-uri,
dimensiunea setului de date crește considerabil.

Clasa \textit{Feature Function} modelează un tip de informație pe care vrem să o extragem din segmentele audio. O instanță primește o funcție care,
pentru un segment audio va returna o informație extrasă din acel segment, parametrii adiționali cu care va fi apelată funcția specificată și
un nume prin care vom identifica informația extrasă. Această abordare asigură elementul de flexibilitate al pipeline-ului. Pentru un tip nou
de informație este necesară doar definirea unei noi funcții.

Informațiile extrase pentru un segment pot fi salvate individual ca JSON sau \textit{Tensors} serializați binar sau pot fi adăugate ca JSON într-un fișier ce
conține informațiile extrase din toate segmentele. Gama largă de reprezentări permite experimentarea cu diferite tipuri de încărcare a datelor și
diferite tipuri de
inputuri în antrenarea rețelei.

În continuare sunt descrise două tipuri de informații care au fost extrase din segmentele audio.

\section{Spectrograme Mel}

\begin{figure}[h]
\includegraphics[width=0.5\textwidth]{melspec.png}
\centering
\captionsetup{justification=centering,font=small,labelfont=small}
\caption{Spectrogramă Mel pentru un segment din setul de date}
\label{fig:f12}
\end{figure}

O spectrogramă este o metodă de vizualizare a evoluției spectrului de frecvențe și intensitatea lor în timp, care se obține aplicând o
transformare Fourier de scurtă durată asupra formei de undă. Transformarea Fourier de scurtă durată se obține aplicând transformări Fourier discrete
pe ferestre de mici dimensiuni ale formei de undă care de obicei se suprapun \cite{spectrograms}. Transformarea Fourier discretă face trecerea de
la domeniul de timp la domeniul de frecvență, descompunând semnalul audio în frecvențele ce îl compun, comparând semnalul cu sinusoidele
diferitelor frecvențe la diferite faze de undă după cum se poate observa în figura \ref{fig:f8}.

\begin{figure}[h]
\includegraphics[width=0.65\textwidth]{ft.png}
\centering
\captionsetup{justification=centering,font=small,labelfont=small}
\caption{Trecerea de la domeniul de timp la domeniul de frecvență (imagine preluată din \textit{Demystifying the Fourier Transform} \protect\footnote[1])}
\label{fig:f8}
\end{figure}
\footnotetext[1]{\url{https://youtu.be/XQ45IgG6rJ4}}

În urma aplicării unei transformări Fourier de scurtă durată obținem o evoluție în timp a frecvențelor prezente în semnalul audio.
Atât amplitudinea cât și frecvențele rezultate sunt reprezentate liniar în urma aplicării acestei transformări. Cu toate acestea, modul în care
oamenii percep aceste valori este logaritmic. De aceea, pentru a apropia această reprezentare de modul în care sunetele sunt percepute,
amplitudinea este convertită în decibeli, iar frecvența este convertită conform scării Mel $m = 2595 \cdot \log (1+\frac{f}{500})$, unde $f$ este
frecvența exprimată în Hertz și $m$ este frecvența exprimată în Mel. Se plasează pe intervalul dintre cea
mai joasă frecvență Mel și cea mai înaltă frecvență Mel un număr de puncte echidistante egal cu numărul de filtre mel dorite (128) care vor fi
centrul filtrelor triunghiulare Mel \ref{fig:f9}.

\begin{figure}[h]
\includegraphics[width=0.55\textwidth]{melbancs.png}
\centering
\captionsetup{justification=centering,font=small,labelfont=small}
\caption{Filtre Mel (\url{http://siggigue.github.io/pyfilterbank/melbank.html})}
\label{fig:f9}
\end{figure}

Pentru a obține spectrograma
Mel, spectrograma rezultată în urma transformării Fourier de scurtă durată cu amplitudinea exprimată în decibeli este înmulțită cu filtrele Mel
calculate.

Această reprezentare a fost aleasă deoarece este apropiată din punct de vedere perceptual de modul în care urechea umană procesează
sunetele, și are rezultate mai bune în literatură în problema clasificării genurilor muzicale\cite{resut1}, sau clasificării
sunetelor în general\cite{large} decât spectrogramele simple sau MFCC atunci când rețeaua este una convoluțională.

\section{MFCC}

\begin{figure}[h]
\includegraphics[width=\textwidth]{mfcc.png}
\centering
\caption{Coeficienți MFCC pentru un segment din setul de date}
\captionsetup{justification=centering,font=small,labelfont=small}
\label{fig:f11}
\end{figure}

Extragerea coeficienților MFCC \ref{fig:f11} (Mel-Frequency cepstral coefficients) se obține realizând aceiași pași descriși în extragerea unei spectrograme,
o transformare Fourier de discretă, convertirea amplitudinii în decibeli și aplicarea bancului de filtre Mel,
urmați de logaritmarea spectrului Mel, aplicarea unei transformări cosinusoidale discrete pentru
fiecare fereastră a semnalului audio și calcularea derivatelor de ordinul întâi si ordinul al doilea, numite coeficienți delta,
respectiv coeficienți delta delta\cite{mfccs}.

Transformarea cosinusoidală discretă este un proces similar cu transformare Fourier discretă, diferența fiind faptul că transformarea cosinusoidală
discretă returnează valori din mulțimea numerelor reale în timp ce transformarea Fourier discretă poate returna valori complexe.
Astfel, se calculează similaritatea
dintre spectrul Mel logaritmic și mai multe funcții cosinusoidale ce corespund diferitelor frecvențe \ref{fig:f10},
această similaritate rezultată reprezentând un coeficient MFCC.
\pagebreak
\begin{figure}[h]
\includegraphics[width=0.9\textwidth]{dct.png}
\centering
\captionsetup{justification=centering,font=small,labelfont=small}
\caption{Transformarea cosinusoidală discretă (imagine preluată din \url{https://youtu.be/4_SH2nfbQZ8})}
\label{fig:f10}
\end{figure}

În timp ce coeficienții MFCC oferă informații relevante despre formanții și fonemele prezente în semnalul sonor pentru recunoașterea
timbrului și \textit{speech recognition} în general, coeficienții delta și delta delta oferă informații temporale în legătură cu ritmul vorbirii.
Datorită acestor aspecte, coeficienți MFCC sunt frecvenți utilizați în problemele de \textit{speech recognition}, principalul avantaj fiind dimensiunea
redusă în comparație cu spectrogramele Mel. Astfel, în funcție de arhitectura rețelei, a fost preferată utilizarea coeficienților MFCC în favoarea
spectrogramei Mel tocmai datorită dimensiunii reduse, în special în cazurile în care nu s-au utilizat straturi convoluționale.
